% =========================================================
% REEMPLAZO de la subsección \label{sec:val_verificador} en
% 3_2_validacion.tex. Sustituir el bloque completo, desde
% \subsection{Verificación de sesiones...} hasta el comentario
% "% PENDIENTE: incorporar el resultado de la prueba de mutación."
% inclusive.
% =========================================================

\subsection{Verificación de sesiones y validación de la herramienta}
\label{sec:val_verificador}

La herramienta \texttt{verificar\_sesion.py} recorre una sesión ya
guardada y evalúa diez condiciones anómalas: ausencia de archivos de
muestras (\texttt{npy\_ausente}), imposibilidad de reconstruir captura
alguna (\texttt{sin\_reconstruibles}), saturación de la señal
(\texttt{saturacion}), repetición de un mismo contenido binario bajo
la misma escala o bajo escalas distintas
(\texttt{buffer\_repetido\_misma\_escala} y
\texttt{buffer\_repetido\_escala\_distinta}), coexistencia de escalas
diferentes en una misma serie (\texttt{mezcla\_escalas}), deriva del
tiempo de arribo (\texttt{deriva\_onset}), estancamiento de la lectura
de temperatura (\texttt{temperatura\_estancada}), presencia de
capturas marcadas con error (\texttt{error\_flag}) y dispersión
excesiva de la amplitud pico a pico (\texttt{dispersion\_vpp}). Es la
herramienta con la que se detectaron las anomalías descritas en la
\cref{sec:falla_buffer} y en la \cref{sec:falla_regresion}.

Una herramienta de auditoría que nunca se ha visto fallar no es una
herramienta verificada, sino una herramienta afortunada. Su validación
se realizó en dos etapas.

La primera consistió en fabricar, mediante
\texttt{generar\_sesion\_prueba.py}, once sesiones sintéticas: diez
diseñadas para provocar cada una de las condiciones anteriores, y una
undécima deliberadamente limpia. Esta última no es un caso menor. Un
verificador que declarara anomalías en toda sesión que examina
detectaría las diez condiciones sin dificultad y sería inútil; la
sesión limpia es la que exige que no las declare cuando no las hay. El
verificador identifica las diez condiciones y no emite alerta alguna
sobre la sesión limpia.

La segunda etapa responde a una objeción que la primera deja abierta.
El corredor compara las alertas obtenidas contra un conjunto esperado,
y ese conjunto podría haberse derivado de la salida observada del
verificador en lugar del comportamiento pretendido; de ser así, la
prueba pasaría igual con el verificador averiado. Para descartarlo se
realizó una prueba de mutación: se desactivó la emisión de cada una de
las diez condiciones, una a la vez, restaurando el archivo entre
mutaciones, y se registró qué casos dejaban de detectarse.

Las diez mutaciones provocaron la caída del caso correspondiente.
Ninguna sobrevivió, y en ningún caso el conjunto completo quedó en
verde con un chequeo desactivado. Dos de ellas arrastraron además un
caso ajeno: al desactivar \texttt{npy\_ausente} cae también
\texttt{sin\_reconstruibles}, y al desactivar \texttt{mezcla\_escalas}
cae \texttt{buffer\_repetido\_escala\_distinta}. Ambas muertes
cruzadas corresponden a relaciones de implicación entre las
condiciones --- una sesión sin capturas reconstruibles es
necesariamente una sesión a la que le faltan archivos de muestras, y
la repetición de contenido bajo escalas distintas presupone escalas
distintas ---, previstas por análisis del código antes de la prueba y
confirmadas empíricamente por ella.

La mutación dejó un hallazgo que no se buscaba. Al desactivar
\texttt{sin\_reconstruibles}, el código de salida del verificador
permanece en 1: la instrucción de retorno anticipada devuelve una lista
que \texttt{npy\_ausente} ya había poblado en la condición anterior. Un
corredor que juzgara el resultado por el código de salida habría dado
ese caso por aprobado con el chequeo apagado. Lo que impidió el falso
aprobado fue la decisión de comparar el conjunto de alertas emitidas
contra el esperado, y no únicamente el código de salida. La
observación no es anecdótica: reaparece en
\cref{sec:val_advertencias}, donde una comprobación distinta calcula su
código de retorno antes del paso que pretende verificar. En ambos
casos, el código de salida resultó ser el indicador menos informativo
disponible.

Queda una limitación que conviene declarar. La prueba de mutación se
aplicó manualmente y no existe en el repositorio como procedimiento
automatizado: si \texttt{verificar\_sesion.py} incorpora una condición
nueva o reescribe una existente, el resultado anterior deja de ser
válido y nada lo advierte. La herramienta que somete al sistema a
verificación continua está, ella misma, verificada por una prueba de
un solo momento.
